\documentclass[../main.tex]{subfiles}
\begin{document}

\section{Depth bands}
\label{sec:depth}

\begin{claim}{STRONG}{F08}
If the backdoor is assembled in a specific band of blocks, a perturbation aimed at that band should beat the same perturbation spread over all \VitBlocks{}, and the best band should be the same for both placement families. The prediction was that a middle band wins everywhere. It wins for residual dropout and loses for the input-side token mask, which is a result about the placement rather than about depth.
\end{claim}

\Cref{tab:depth-bands} reads every band placement the basis carries at the matched rule, with the paired delta against its family's all-blocks placement. Restricting the recommended token mask to blocks \BandMiddle{} costs \BandFiveEightMinusAllInputSide{} with interval [\BandFiveEightMinusAllInputSideLow{}, \BandFiveEightMinusAllInputSideHigh{}], and to blocks \BandLate{} costs \BandNineOneTwoMinusAllInputSide{} with interval [\BandNineOneTwoMinusAllInputSideLow{}, \BandNineOneTwoMinusAllInputSideHigh{}]. Restricting residual dropout to blocks \BandMiddle{} gains \BandFiveEightMinusAllResidual{} with interval [\BandFiveEightMinusAllResidualLow{}, \BandFiveEightMinusAllResidualHigh{}], to blocks \BandLate{} gains \BandNineOneTwoMinusAllResidual{} with an interval that includes $0$, and to blocks \BandEarly{} costs \BandOneFourMinusAllResidual{} with interval [\BandOneFourMinusAllResidualLow{}, \BandOneFourMinusAllResidualHigh{}]. Banding helps the residual-adjacent family and hurts the input-side family, so the earlier statement that a band beats the full stack is a property of residual dropout and not of depth.

% generated by scripts/paper/tab_depth_bands.py from results/coverage/coverage.json, results/<folder>/psbd_metrics.json (65 cells) at 2026-09-11T04:12:43.663538+00:00, commit 4a04e5dd204b14d8368fb7eeef7092bd331e00d9-dirty. Do not edit by hand.
\begin{table}[htbp]
\centering
\small
\caption{Depth-band placements at the matched 0.6 rate, mean AUROC at q0.25, the paired delta against the family's all-blocks placement with a 5000-resample bootstrap 95\% interval, and the mean achieved clean-validation shift ratio at the chosen rate. Token masking at the attention input restricted to blocks 1 to 4 has not been swept yet on any model in the panel of 65 backdoored models and prints pending throughout.}
\label{tab:depth-bands}
\begin{adjustbox}{max width=\linewidth}
\begin{tabular}{lrrrlr}
\toprule
placement & n & mean AUROC matched06 & delta vs all blocks & 95\% CI & mean achieved shift \\
\midrule
input-side, all blocks & 65 & 0.926 & -- & -- & 0.595 \\
input-side, blocks 5-8 & 65 & 0.863 & -0.063 & [-0.100, -0.031] & 0.603 \\
input-side, blocks 9-12 & 65 & 0.864 & -0.062 & [-0.108, -0.021] & 0.489 \\
input-side, blocks 1-4 & \pending & \pending & \pending & \pending & \pending \\
residual-adjacent, all blocks & 65 & 0.841 & -- & -- & 0.658 \\
residual-adjacent, blocks 1-4 & 65 & 0.779 & -0.062 & [-0.081, -0.042] & 0.582 \\
residual-adjacent, blocks 5-8 & 65 & 0.893 & +0.052 & [+0.027, +0.079] & 0.565 \\
residual-adjacent, blocks 9-12 & 65 & 0.874 & +0.033 & [-0.012, +0.076] & 0.598 \\
\bottomrule
\end{tabular}
\end{adjustbox}
\end{table}


The halves of the table read differently for a mechanistic reason. Residual dropout at all blocks masks the stream \VitBlocks{} times and the adaptive rule has to pick a small rate to stay inside its range, so removing the early blocks from the mask removes disturbance before the trigger has been routed anywhere, where there is nothing to disrupt, and concentrates the same clean-validation shift on the blocks where \cref{sec:latent} finds the routing. The token mask at the attention input perturbs what each attention block reads and leaves the stream itself untouched, so a band of it cannot drive the clean-validation shift ratio as high as the full stack can and the band rows sit at lower achieved shift in the table's last column, which is the strength confound reappearing inside a single placement. The early band of the input-side family, blocks \BandEarly{}, has no cache anywhere in the panel and prints \pending{}. It is the single row that would separate the routing account from the strength account, since routing predicts it to be the worst band of its family and strength predicts it to be indistinguishable from the other bands, and it is listed in \cref{sec:planned} with its command.

An earlier rule of this project selected the band from the layer at which an attack's direction reaches the class token, and that rule was pre-registered against a pair of held-out attacks and failed on both. The band effect is measured. Its selection rule is not, and the table's residual rows are the evidence for a fixed middle band rather than for a per-attack choice, which a defender who does not know the attack could not make anyway.

\end{document}
